
\subsection{A-Levels}

A-levels are the terminal, standardized qualifications taken at the end of upper-secondary education (age 16–18) and form the principal basis for university admission (Figure~1). Applicants may apply to up to five courses via UCAS, typically by a January deadline; because the examinations are held after applications close, applications are supported by teacher-predicted grades. Universities then issue conditional offers (usually by May) that specify the grades required for admission. Applicants nominate a firm and an insurance choice; these choices establish the binding order in which offers are accepted.

Final placement is determined by the centrally awarded A-level grades. A-level scripts are marked by examination boards: markers are blind to candidates' identities and to application information, and letter grades are assigned according to standardised grade boundaries. Once actual results are released in August, UCAS assigns students according to the ranking they submitted and the conditions of the conditional offers (the firm choice is accepted if its conditions are met; otherwise the insurance choice may be used; failing both, the student can enter Clearing). Because grading and placement follow this centralised, rule-based procedure, teachers and schools have no formal discretion over the final marks or over the contractual placement outcome.

A-level scripts are marked centrally by examination boards; markers are blind to candidates' identities and demographic attributes. Letter grades are awarded according to grade boundaries that examiners set annually. Those boundaries are calibrated with reference to historical performance and other standardisation procedures so that the difficulty of achieving a given grade is broadly comparable across years. We confirm this property empirically by predicting A-level attainment from pre-treatment observables (here, GCSE mathematics and English) and comparing the conditional expectation of A-level outcomes across non-COVID years (\Cref{fig:event_study}): visual inspection of binned conditional expectation functions, pooled regressions with year × predicted-score interactions, and formal equivalence tests within narrow predicted-score bins reveal no meaningful shifts in the conditional means. Complementary checks — subject-by-exam-board splits — produce the same conclusion, supporting the claim that, conditional on prior attainment, A-level outcomes are stable across non-COVID cohorts.

Subject choice and course provision vary across the system. There are well over 100 distinct A-level subjects, and schools and sixth-form colleges decide which subjects to offer according to local demand and resources. Some pupils remain at their secondary school for sixth-form study, while others transfer to separate sixth-form or further-education colleges; this spatial and institutional dispersion affects which subjects pupils can take. Universities often specify required subjects for particular courses, and highly selective institutions tend to favour so-called ``facilitating'' subjects; independent schools frequently encourage uptake of these subjects and invest in resources (e.g. specialised equipment and teaching) to support them.

\subsection{University application system in the UK}
The UK uses a centralised admissions system (UCAS) to match applicants to university courses. Applicants may apply to up to five courses through UCAS; universities reply—typically before students sit their A-levels—with acceptances, rejections, or conditional offers that specify the grades required for admission. After receiving offers, applicants must choose two to keep: a \emph{firm} choice and an \emph{insurance} choice, and withdraw any remaining offers. These ranked choices determine allocation once actual A-level results are released: if the firm offer’s conditions are met the firm place is confirmed, otherwise the insurance offer may be used; failing both, the applicant may enter Clearing to seek an alternative place.

\subsection{The British Education Sector for Higher Education}
The UK has on average 3,300 university programs per year, matching with on average 570,000 students per a year. 3,500 secondary schools provide students with education to prepare for exams.

British secondary schools are mostly run publicly. 93\% of students participate in public schools and only 7\% in private schools ((\Cref{tab:balance})). Private schools, equally known as independent schools, run mainly on tuition collected from their pupils' parents. Past literature documents that attendants of independent schools are often children with rich parents. Henseke et al (2021) finds that 92 percent of the students attending independent schools belong to households in the top 10th decile group of the income distribution. 90\% of the students in independent schools proceed to higher education after graduation, with 4\% attending Oxford of Cambridge and 52\% attending one of the top 25th ranked universities. Students from independent schools apply to better universities more than students from public schools. Almost 85\% of the students apply to at least one university of the Russell group, while 45\% of the students from public schools apply. Independent school students apply less to lower-tier schools than students from public schools. Almost 40\% of the students apply to lower-tier schools, while 80\% of the students from public schools apply to at least one lower tier school. Independent school students receive more offer of admission from Russell Group schools than public schools. 80\% of students from independent schools receive at least one offer from a Russell Group school. 

There are multiple types of public school in the UK higher education sector (\Cref{tab:balance}). The most common type of schools are academy schools that consists of 35\% of the total number of schools. Academy schools provide education to pupils living within their local authority jurisdiction. As an exception, grammar schools select students based on their performance on a previous national exam (\emph{year 11}). 

\subsection{COVID-19 pandemic}
The COVID-19 pandemic led the UK authorities to cancel in-person A-level examinations and substitute centrally-marked exams with teacher-based assessments. For the 2020 cohort this change was announced on 18 March 2020, and for the 2021 cohort on 6 January 2021. In place of externally marked scripts, teachers were asked to submit centre-assessed grades for their pupils, drawing on their knowledge of each student and the student’s prior classroom performance and exam history. Exam boards instructed centres to review and, where necessary, lower implausibly high submissions before finalising their returns; nonetheless, the grades submitted by schools frequently stood as the final awards and were used in place of the usual centrally-marked A-level marks.


Crucially, both announcements came after applications had been submitted and universities had issued responses (acceptances, rejections, conditional offers or unconditional offers). The teacher-based grades therefore stood in for the usual, centrally-awarded A-level marks used to satisfy conditional offers and determine final placements. At that point universities had little practical means to control the incoming cohort — their main lever was to tighten conditional offers (raise grade requirements) but many institutions were reluctant to do so for fear of deterring applicants during the pandemic. Importantly, students had already nominated a firm and an insurance choice before teachers assigned centre-assessed grades, so final placements followed the pre-submitted ranking together with the teacher-submitted grades rather than any subsequent adjustment by schools or universities.



\subsection{Consequence of the method change}
The immediate consequence of the method change was substantial grade inflation (Figure~\ref{fig:grade_histogram}). The share of students awarded A* or A rose by roughly 12 percentage points in 2020 and by about 20 percentage points in 2021 relative to pre-pandemic cohorts. Importantly, the shift was not confined to the top tail: the entire grade distribution moved (see Figure~\ref{fig:grade_histogram}), with noticeable changes across middle and lower bands (C–E) as well as at A/A*. In short, teacher-based assessment redistributed mass across multiple grade bands rather than merely increasing top grades.

Those inflated grades translated directly into larger incoming cohorts because universities were obliged to honour conditional offers for students who met the stated requirements. Figures~\ref{fig:CohortSizeGrowthAny} and \ref{fig:CohortSizeGrowthFirm} document this over-acceptance in both 2020 and 2021. The expansion was concentrated among more selective programmes: the most competitive programmes—those with low pre-pandemic acceptance rates—experienced the largest proportional increases in cohort size, with some programmes almost doubling intake despite their typically small pre-pandemic capacities.

Figure~\ref{fig:CohortSizeGrowthFirm} shows that a substantial part of the expansion came from students being placed at their firm choice. In particular, programmes that lay below the 30th percentile of pre-pandemic acceptance rates increased their intake by more than twice the pre-pandemic average, indicating that the grade shifts interacted with the existing selectivity hierarchy to produce outsized growth at the top end of the market.


%Table [] reports the simulation results for income sorting across universities by grading regime.  Column 1, 2, and 3 reports the correlation between university rank and the average parental income of students for incoming student cohort for standardized, non-standardized grading scheme for 2020 and 2021, respectively. The results indicate how the sorting pattern across university by income strengths for both 2020 and 2021 grading patterns. A one percent increase in the selectivity of a university, measured by the incoming cohorts average tariff scores, is associated with a 0.44 increase in the standard deviation of the parental income distribution.  The slope is statistically significantly larger than the magnitude of the coefficient with the standardized exam grading patterns. 

%Table [] reports the simulation results for the average academic competence of the incoming cohort measured by students GCSE math score. The non-standardized regime predicts that the average sorting pattern between students academic competence and the selectivity of the university strengths, with a one percent increase in the universities selectivity leading to a 1.04 standard deviation point reduction in the incoming cohorts average GCSE math score. The estimate for the standardized exam regime was -0.69, and it is statistically smaller in magnitude than the non-standardized regime. [Without endogenous application.] 

%Table [] reports the simulation results for the distribution of females among incoming across universities. Although the female students receive higher grades than male students in the non-standardized testing regime, the regression coefficient between female share and the selectivity of the university decreases. A one-percent increase in the selectivity of the university is associated with a 0.2 percentage point decrease in the share of female students with the non-standardized grading patterns. The estimate is -0.07 percentage point for the standardized grading regime.
